\section{Reproducibility Snapshot}
\label{app:reproducibility}

\begin{table}[h]
  \caption{Frozen knowledge-graph snapshot and contract-level edge coverage.  Coverage is
  structural and is not a correctness rate.}
  \label{tab:dataset-stats}
  \centering
  \small
  \begin{tabular}{llrr}
    \toprule
    Block & Item & Count & Share \\
    \midrule
    \multirow{4}{*}{Graph nodes}
      & Contract/notice & \KGContracts & \\
      & Organisation & \KGOrgs & \\
      & CPV & \KGCPV & \\
      & Evidence & \KGEvidence & \\
    \midrule
    \multirow{4}{*}{Contract coverage}
      & Buyer edge & \KGBuyerCoverage & 99.999\% \\
      & Supplier edge & \KGSupplierCoverage & 94.87\% \\
      & Category edge & \KGCategoryCoverage & 76.52\% \\
      & At least one evidence edge & \KGEvidenceCoverage & 99.99\% \\
    \bottomrule
  \end{tabular}
\end{table}

\subsection{Code and artifact state}

The repository base commit inspected for this draft is
\texttt{d15ec0a3bd5b5f295f9828e44067094b73ab2cd7}.  The working tree also contains uncommitted
correctness patches and staged data artifacts; the commit alone is therefore not an exact
reproduction capsule.  Before submission, the authors must create a clean tagged revision and
record the patch, Python environment lock, CUDA/driver versions, model revisions, and command
lines.  No new model training or batch inference was run to prepare this draft.

\begin{table}[h]
  \caption{Content hashes for the primary frozen inputs.}
  \label{tab:artifact-hashes}
  \centering
  \scriptsize
  \begin{tabular}{p{0.44\linewidth}p{0.48\linewidth}}
    \toprule
    Repository path & SHA-256 \\
    \midrule
    \path{data/qa/pacs_v1/test_channel_a.jsonl}
      & \texttt{1869558ad2764b92888df47b093dfc2f}\newline
        \texttt{f693d2fdfdd6544f1fa3734d017cf2bb} \\
    \path{data/qa/pacs_v1/test_channel_b.jsonl}
      & \texttt{36100c4bfb7b34ef6fe1d72c62d33a62}\newline
        \texttt{d069501aad97c0e28a172345a901339b} \\
    \path{data/training/llamafactory_compose_v3/compose_sft_train.json}
      & \texttt{a80945184aeb9531321b26ef57286656}\newline
        \texttt{8707eeba10b7dc56ec34f8e35f00a20d} \\
    \path{data/training/llamafactory_compose_v3/compose_sft_val.json}
      & \texttt{7cba3d43610a90216e5d7655319178df}\newline
        \texttt{3a66f1728d1133b704808b0a8ba2a899} \\
    \path{data/kg/nodes/contract_nodes.parquet}
      & \texttt{246caad93e1c8b184d3714a280190efb}\newline
        \texttt{3c238f4cab73cd3bf1f36b38de307bf1} \\
    \path{data/kg/nodes/org_nodes.parquet}
      & \texttt{42f9542833275ea7a57feb2cc4f4e64}\newline
        \texttt{a507513f35b2e420a449d938b68b867af} \\
    \bottomrule
  \end{tabular}
\end{table}

Per-item PACS predictions and summaries are stored under
\texttt{data/qa/compose\_probe\_v1/}.  The primary summaries are
\texttt{summary\_pacstest\_base\_a.json},
\texttt{summary\_pacstest\_teacher\_a.json},
\texttt{summary\_pacstest\_v3\_a.json}, and
\texttt{summary\_pacstest\_v3\_b.json}.  The corresponding JSON Lines files retain the raw
model response, parsed tree, execution answer, oracle, outcome, and correctness flag.

The 260-item runtime summaries are under
\texttt{outputs/eval/matrix\_v2/fully-local-qwen/},
\texttt{.../cicada-qwen3-zeroshot/}, and \texttt{.../teacher\_r3/}.  These paths contain an
internal historical label and currently resolve through a machine-specific absolute symlink;
that label is not used as the method's paper name.  A release must replace the absolute link with
content-addressed manifests and publish model, prompt, environment, and evaluator hashes, which
is the condition under which the \RuntimeN{}-item runtime result could be promoted from a frozen
artifact to primary evidence.

\subsection{Pipeline entry points}

A full deterministic data build follows this order:
\begin{verbatim}
.venv/bin/python -B pipelines/01_ingest.py
.venv/bin/python -B pipelines/02_er_phase1.py
.venv/bin/python -B pipelines/03_er_phase2.py
.venv/bin/python -B pipelines/04_extract.py
.venv/bin/python -B pipelines/40_build_kg.py
.venv/bin/python -B pipelines/41_validate_kg.py
\end{verbatim}
Partial year/sample runs require explicit alternate outputs in the current code and must never
be used to overwrite the canonical artifact.  Exact raw-data acquisition commands and licences
remain to be added to the release manifest.

Program supervision is built with \texttt{scripts/build\_compose\_train.py}, exported with the
compose export script, and trained with
\texttt{configs/training/qwen3\_8b\_compose\_sft\_v3\_qlora.yaml}.  The stored training summary
reports train loss 0.0171394, evaluation loss 0.0096943, and runtime 8270.93 seconds.  These losses
are diagnostics only because the exported split is not fully deduplicated.

PACS generation uses, in order,
\texttt{scripts/pacs/generate.py}, \texttt{render\_channels.py},
\texttt{surface\_independent.py}, \texttt{naturalize.py}, and \texttt{seal.py}.  Evaluation uses
\texttt{scripts/run\_compose\_probe\_eval.py}.  A release command must record the serving
endpoint, exact model revision, API date/version, prompt hash, maximum tokens, concurrency,
temperature, guided-decoding flag, reflection budget, and scorer revision.

\subsection{Checklist before final submission}

The final artifact should contain: (i) a clean source tag and dirty-tree check; (ii) raw-data and
KG manifests; (iii) model/adapter hashes or durable release links; (iv) exact prompts and JSON
schemas; (v) strict scorer unit tests covering Boolean versus numeric values; (vi) route,
repair, and failure counts; (vii) random seeds and repeated-run policy; (viii) hardware and
software versions; and (ix) a one-command script that regenerates every table from JSON/Parquet
without copying numbers by hand.
